\section{Conclusioni}
Nel recupero conversazionale di RETECO, la scelta di come costruire la lista dei candidati incide già prima di introdurre calcolo aggiuntivo nel decoder. Nella validazione interna disponibile la fusione contestuale ottiene la migliore media tra le sette configurazioni esaminate, pur cedendo a un singolo retriever in diversi domini. Questo risultato individua due condizioni sperimentali informative: un pool lessicale riproducibile e un pool ibrido più competitivo in media.

L'implementazione consente ora di porre la stessa domanda sulla profondità in due condizioni distinte: riordinare evidenze già reperite oppure formulare una risposta a partire da evidenze fornite. Il passo decisivo per il preprint è misurare, a turno e input fissati, la differenza rispetto a una sola applicazione dello strato, insieme al tempo necessario per ottenerla. Da quel confronto dipenderà la conclusione sull'utilità della ricorrenza.
